\centerline{\bf \Huge ОТА. Делимость I}

\begin{flushright}
        \scriptsize{Паника — признак слабости духа\\Пешши из ДжоДжо}
\end{flushright}

\textbf{Определение 1.} Натуральное число больше 1 называется простым, если оно делится нацело только на единицу и само себя

\textbf{Определение 2.} Натуральное число называется составным, если оно делится нацело не только на единицу и само себя.

\textbf{Определение 3.} Числа называются взаимно простыми, если их наибольший общий делитель равен 1.

Основная теорема арифметики. Любое число больше 1 может быть разложено в произведение одного или больше простых множителей, причем это разложение единственно.

$$
n=p_1^{\alpha_1} p_2^{\alpha_2} p_3^{\alpha_3} \cdots p_k^{\alpha_k}
$$

\textbf{Важная лемма!}
Докажите, что если $a b$ делится на простое число $p$, то одно из чисел $a$ и $b$ делится на $p$

\problem{1}
Докажите, что произведение любых трёх последовательных натуральных чисел делится на 6.

\problem{2}
Про $a, b$ известно, что ни одно из них не кончается на ноль, а произведение 10000 . Какие значения может принимать их сумма?

\problem{3}
Существуют ли 3 натуральных числа таких, что ни одно из них не делится на другое, произведение любых двух из них делится на третье?

\problem{4}
Существует ли n такое, что $n!$, что нулей на конце ровно $5?$

\problem{5}
Если в выражение $\mathrm{n}^2-\mathrm{n}+41$ подставлять числа $\mathrm{n}=1,2,3,4,5$, то получатся простые числа $41,43,47,53,61$. Верно ли, что при любом n получится простое число?

\problem{6}
Докажите, что у составного числа $a$ найдется такой простой делитель $p$, что $p^2 \leqslant a$.